
\begin{table}[H]
\centering
\caption{S\&P500 low- and high-volatility bucket diagnostics}
\label{tab:sp500-lowvol-buckets}
\scriptsize
\resizebox{\textwidth}{!}{%
\begin{tabular}{llrrrr}
\toprule
Model & Bucket & Mean truth & Mean pred. & Bias & MSE \\
\midrule
\model{HAR_M} & Q1 & 15.92 & 16.72 & 0.81 & 3.39 \\
\model{HAR_M} & Q5 & 54.87 & 53.35 & -1.52 & 32.59 \\
\model{GHAR_M_IV} & Q1 & 15.92 & 16.68 & 0.76 & 3.28 \\
\model{GHAR_M_IV} & Q5 & 54.87 & 53.89 & -0.98 & 30.81 \\
\model{GNNHAR1L_M} & Q1 & 15.92 & 29.98 & 14.06 & 203.72 \\
\model{GNNHAR1L_M} & Q5 & 54.87 & 53.53 & -1.34 & 31.26 \\
\model{GNNHAR1L_M_IV} & Q1 & 15.92 & 29.98 & 14.06 & 203.55 \\
\model{GNNHAR1L_M_IV} & Q5 & 54.87 & 53.84 & -1.03 & 30.88 \\
\model{GNNHAR2L_Q_IV} & Q1 & 15.92 & 26.97 & 11.05 & 127.92 \\
\model{GNNHAR2L_Q_IV} & Q5 & 54.87 & 53.61 & -1.27 & 507.35 \\

\bottomrule
\end{tabular}
}
\begin{minipage}{0.96\textwidth}
\footnotesize
\emph{Notes.} Buckets are quintiles of the realized-volatility proxy.  Q1 is the lowest-volatility quintile and Q5 is the highest-volatility quintile.
\end{minipage}
\end{table}
